% GENERATED FILE. Edit canonical lessons, then run npm run generate:books.
% canonical-source-hash: fnv1a64:fc1ea1b573b1b404
% canonical-lessons: TA-C195-kontu-va, TA-C195-tiruppik-kotu, TA-C195-ottu2, TA-C195-payanpatuttu, TA-C195-muyarci-cey

\chapter{To bring, To give back, To drive, To use, To try}
\label{ch:ta-to-bring-to-give-back-to-drive-to-use-to-try}
\hlchaptermodality{195}

\begin{chapteropening}
\textbf{By the end of this chapter:} \emph{I can say to bring, to give back, to drive, to use and to try.}

Complete the last lesson of chapter 195: I can say to bring, to give back, to drive, to use and to try.
\end{chapteropening}

\section[\texorpdfstring{\ta{கொண்டு} \ta{வா} (koṇṭu vā)}{கொண்டு வா (koṇṭu vā)}]{\ta{கொண்டு} \ta{வா} (koṇṭu vā) — to bring}

\label{lesson:TA-C195-kontu-va}



\begin{quote}
Before the new one: say the Tamil for to deep-fry, then the Tamil for to knead.
\end{quote}

\subsection*{You'll want to know: \ta{கொண்டு} \ta{வா}}
\textbf{\ta{கொண்டு} \ta{வா}} — \emph{koṇṭu vā} — \textquotedblleft{}to bring\textquotedblright{}.

\begin{grammarlens}[title={What you've built}]
One more word for this chapter.
\end{grammarlens}

\subsection*{Your turn}
\begin{itemize}
\raggedright
  \item \emph{Say it:} \emph{koṇṭu vā}
  \item \emph{Say it:} \emph{koṇṭu vā}, once more
  \item \emph{Say it:} say \emph{koṇṭu vā}
  \item \emph{Recall:} say the Tamil for to deep-fry, then the Tamil for to knead, then say \emph{koṇṭu vā} again
\end{itemize}

\begin{tcolorbox}[breakable,colback=teal!4,colframe=teal!35!black,title={Before you move on}]
What does \ta{கொண்டு} \ta{வா} mean? (\textquotedblleft{}To bring\textquotedblright{}.) Say it once more.
\end{tcolorbox}
\section[\texorpdfstring{\ta{திருப்பிக்} \ta{கொடு} (tiruppik koṭu)}{திருப்பிக் கொடு (tiruppik koṭu)}]{\ta{திருப்பிக்} \ta{கொடு} (tiruppik koṭu) — to give back}

\label{lesson:TA-C195-tiruppik-kotu}



\begin{quote}
Before the new one: say the Tamil for to knead, then the Tamil for to bring.
\end{quote}

\subsection*{You'll want to know: \ta{திருப்பிக்} \ta{கொடு}}
\textbf{\ta{திருப்பிக்} \ta{கொடு}} — \emph{tiruppik koṭu} — \textquotedblleft{}to give back\textquotedblright{}.

\begin{grammarlens}[title={What you've built}]
One more word for this chapter.
\end{grammarlens}

\subsection*{Your turn}
\begin{itemize}
\raggedright
  \item \emph{Say it:} \emph{tiruppik koṭu}
  \item \emph{Say it:} \emph{tiruppik koṭu}, once more
  \item \emph{Say it:} say \emph{tiruppik koṭu}
  \item \emph{Recall:} say the Tamil for to knead, then the Tamil for to bring, then say \emph{tiruppik koṭu} again
\end{itemize}

\begin{tcolorbox}[breakable,colback=teal!4,colframe=teal!35!black,title={Before you move on}]
What does \ta{திருப்பிக்} \ta{கொடு} mean? (\textquotedblleft{}To give back\textquotedblright{}.) Say it once more.
\end{tcolorbox}
\section[\texorpdfstring{\ta{ஓட்டு} (ōṭṭu)}{ஓட்டு (ōṭṭu)}]{\ta{ஓட்டு} (ōṭṭu) — to drive}

\label{lesson:TA-C195-ottu2}



\begin{quote}
Before the new one: say the Tamil for to bring, then the Tamil for to give back.
\end{quote}

\subsection*{You'll want to know: \ta{ஓட்டு}}
\textbf{\ta{ஓட்டு}} — \emph{ōṭṭu} — \textquotedblleft{}to drive\textquotedblright{}.

\begin{grammarlens}[title={What you've built}]
One more word for this chapter.
\end{grammarlens}

\subsection*{Your turn}
\begin{itemize}
\raggedright
  \item \emph{Say it:} \emph{ōṭṭu}
  \item \emph{Say it:} \emph{ōṭṭu}, once more
  \item \emph{Say it:} say \emph{ōṭṭu}
  \item \emph{Recall:} say the Tamil for to bring, then the Tamil for to give back, then say \emph{ōṭṭu} again
\end{itemize}

\begin{tcolorbox}[breakable,colback=teal!4,colframe=teal!35!black,title={Before you move on}]
What does \ta{ஓட்டு} mean? (\textquotedblleft{}To drive\textquotedblright{}.) Say it once more.
\end{tcolorbox}
\section[\texorpdfstring{\ta{பயன்படுத்து} (payaṉpaṭuttu)}{பயன்படுத்து (payaṉpaṭuttu)}]{\ta{பயன்படுத்து} (payaṉpaṭuttu) — to use}

\label{lesson:TA-C195-payanpatuttu}



\begin{quote}
Before the new one: say the Tamil for to give back, then the Tamil for to drive.
\end{quote}

\subsection*{You'll want to know: \ta{பயன்படுத்து}}
\textbf{\ta{பயன்படுத்து}} — \emph{payaṉpaṭuttu} — \textquotedblleft{}to use\textquotedblright{}.

\begin{grammarlens}[title={What you've built}]
One more word for this chapter.
\end{grammarlens}

\subsection*{Your turn}
\begin{itemize}
\raggedright
  \item \emph{Say it:} \emph{payaṉpaṭuttu}
  \item \emph{Say it:} \emph{payaṉpaṭuttu}, once more
  \item \emph{Say it:} say \emph{payaṉpaṭuttu}
  \item \emph{Recall:} say the Tamil for to give back, then the Tamil for to drive, then say \emph{payaṉpaṭuttu} again
\end{itemize}

\begin{tcolorbox}[breakable,colback=teal!4,colframe=teal!35!black,title={Before you move on}]
What does \ta{பயன்படுத்து} mean? (\textquotedblleft{}To use\textquotedblright{}.) Say it once more.
\end{tcolorbox}
\section[\texorpdfstring{\ta{முயற்சி} \ta{செய்} (muyaṟci cey)}{முயற்சி செய் (muyaṟci cey)}]{\ta{முயற்சி} \ta{செய்} (muyaṟci cey) — to try}

\label{lesson:TA-C195-muyarci-cey}



\begin{quote}
Before the new one: say the Tamil for to drive, then the Tamil for to use.
\end{quote}

\subsection*{You'll want to know: \ta{முயற்சி} \ta{செய்}}
\textbf{\ta{முயற்சி} \ta{செய்}} — \emph{muyaṟci cey} — \textquotedblleft{}to try\textquotedblright{}.

\begin{grammarlens}[title={What you've built}]
One more word for this chapter.
\end{grammarlens}

\subsection*{Your turn}
\begin{itemize}
\raggedright
  \item \emph{Say it:} \emph{muyaṟci cey}
  \item \emph{Say it:} \emph{muyaṟci cey}, once more
  \item \emph{Say it:} say \emph{muyaṟci cey}
  \item \emph{Recall:} say the Tamil for to drive, then the Tamil for to use, then say \emph{muyaṟci cey} again
\end{itemize}

\begin{tcolorbox}[breakable,colback=teal!4,colframe=teal!35!black,title={Before you move on}]
What does \ta{முயற்சி} \ta{செய்} mean? (\textquotedblleft{}To try\textquotedblright{}.) Say it once more.
\end{tcolorbox}
